\documentclass[aspectratio=169]{beamer}
\usetheme{Madrid}
\usecolortheme{default}
\title{Planificación y Operación de Sistemas Eléctricos de Potencia}
\subtitle{Mapa general del curso y ruta de aprendizaje}
\author{Prof. Wilian Guamán Cuenca}
\date{2026}
\begin{document}
\begin{frame}\titlepage\end{frame}

\begin{frame}{Propósito de la asignatura}
La asignatura estudia cómo formular, implementar y analizar modelos de operación y planificación de sistemas eléctricos de potencia.
\vspace{0.3cm}
\begin{itemize}
\item Convertir decisiones técnicas en modelos matemáticos resolubles.
\item Implementar modelos con AMPL, Excel y Python.
\item Interpretar resultados desde el punto de vista eléctrico, operativo y económico.
\end{itemize}
\end{frame}

\begin{frame}{Horizonte temporal de los problemas eléctricos}
\begin{tabular}{lll}
\textbf{Horizonte} & \textbf{Problema típico} & \textbf{Herramienta}\\
Segundos--minutos & Control y estabilidad & Simulación dinámica\\
Horas--días & Despacho, UC, OPF & LP, QP, MILP, NLP\\
Semanas--meses & Hidrología, mantenimiento & Escenarios\\
Años & GEP/TNEP & MILP multietapa\\
Décadas & Transición energética & Robustez y sensibilidad\\
\end{tabular}
\end{frame}

\begin{frame}{Ruta de aprendizaje}
\begin{enumerate}
\item Optimización matemática aplicada.
\item AMPL como lenguaje algebraico.
\item Operación de corto plazo.
\item Flujo óptimo de potencia.
\item Proyección de demanda.
\item Expansión de transmisión.
\item Expansión de generación.
\end{enumerate}
\end{frame}

\begin{frame}{Flujo de herramientas}
\centering
Excel/CSV $\rightarrow$ Python $\rightarrow$ AMPL $\rightarrow$ Solver $\rightarrow$ Resultados $\rightarrow$ Python/Excel
\vspace{0.4cm}

\begin{itemize}
\item Excel organiza datos de entrada.
\item Python valida, transforma y visualiza.
\item AMPL formula y resuelve modelos algebraicos.
\item Los resultados se interpretan con tablas, gráficos y sensibilidad.
\end{itemize}
\end{frame}
\end{document}
